\section{A maximal-isotropic boundary for real rank-four inducing groups}
\label{sec:rank-four-real-boundary}

\begin{theorem}[D-BI-12: removal of the real rank-four cases]
\label{D-BI-12}
In the situation of D-BI-11, the reflection normal closure $R$ of
the four-dimensional primitive inducing image cannot be a real
reflection group.  This conclusion allows a larger inducing field
$L\supsetneq K$: the maximal isotropic space constructed in the
proof is defined over the original quadratic field $K$.

Consequently, among the eight non-Eisenstein class-number-one fields,
the remaining rank-four cases of D-BI-11 have
\[
\begin{array}{c|c|c}
 K&R&L\\\hline
 \mathbf Q(i)&G_{29},G_{31}&\mathbf Q(\zeta_8)\text{ or }
                                      \mathbf Q(\zeta_{12})\\
 \mathbf Q(i)&G_{32}&\mathbf Q(\zeta_{12})\\
 \mathbf Q(\sqrt{-2})&G_{29},G_{31}&\mathbf Q(\zeta_8).
\end{array}                                                    \tag{BI.31}
\]
All rows still require the twelve-branch profile and field descent
in D-BI-11.  No existence or dominance assertion is made for them.
\end{theorem}
\begin{proof}
Suppose $R$ is real.  If $J=R$, its character field is real, which
contradicts $K\subset\mathbf Q(\chi_\tau)$.  The remaining branch
pattern in D-BI-11 therefore consists of ten reflections and two
rank-two outside inertias.  Hurwitz moves place the latter at the
end of the tuple.  Moving an outside element past a reflection
conjugates that reflection by a normalizer of $R$, so all ten
reflection entries stay real reflections in $R$.  We may work after
forgetting or reordering the branch labels; these are finite changes
and leave the period rank unchanged.

Choose real orthogonal coordinates for $R$.  Every complex unitary
normalizer is scalar times real orthogonal, by Cohen's Lemma~2.12
and uniqueness of the positive invariant form.  Write the two
outside inertia matrices as
\[
                   g_1=\lambda A_1,\qquad
                   g_2=\lambda^{-1}A_2,                      \tag{BI.32}
\]
where $A_1,A_2$ are real orthogonal and $\lambda$ is nonreal of
absolute value one.  To justify the choices, the product $g_1g_2$
belongs to $R$, so the product of their two unit scalar phases is
real, hence $1$ or $-1$; absorb the latter sign into $A_2$.
If both phases were real, every inertia and hence $J$ would have
real character, again contradicting $K\subset L$.  Thus both
phases are nonreal and (BI.32) is valid.

Put $w=(g_1g_2)^{-1}\in R$.  The ten reflections give a
factorization of $w$.  Let $\ell$ be its shortest reflection length
in the finite real reflection group $R$.  Lewis--Reiner's published
pair-extraction result applies to an arbitrary element, not only
to a Coxeter element:
\cite[Corollary~1.4, p.~1459; proof \S4.3, pp.~1480--1481]{BI:LR}.
It changes this length-ten factorization by Hurwitz moves into
\[
 (r_1,r_1,\ldots,r_q,r_q,s_1,\ldots,s_\ell),
 \qquad q=\frac{10-\ell}{2},\qquad s_1\cdots s_\ell=w,
                                                               \tag{BI.33}
\]
where the final factorization is shortest.  Carter's criterion,
in the precise real-root-system form stated in
\cite[Proposition~4.1, p.~1473]{BI:LR}, says that the roots of
$s_1,\ldots,s_\ell$ are linearly independent.  In particular
$\ell\le4$.

We record explicitly the fixed-space consequence.  If those roots
are $\alpha_1,\ldots,\alpha_\ell$, the telescoping identity
\[
 s_1\cdots s_\ell-1
       =\sum_{j=1}^{\ell}s_1\cdots s_{j-1}(s_j-1)
\]
has image terms along the roots
$\beta_j=s_1\cdots s_{j-1}\alpha_j$.
Each $\beta_j$ is $\alpha_j$ plus a combination of earlier
$\alpha_i$, so the $\beta_j$ are independent.  Applying the
identity to a fixed vector of $w$ shows that every coefficient,
and hence every $(s_j-1)x$, is zero.  Therefore
\[
 \ker(w-1)=\bigcap_{j=1}^{\ell}\ker(s_j-1),\qquad
                \dim\ker(w-1)=4-\ell.                      \tag{BI.34}
\]
The parity of the original ten-reflection factorization gives
$\ell\in\{0,2,4\}$ and $q\in\{5,4,3\}$.

Let $H_0\le J$ be generated by
$s_1,\ldots,s_\ell,g_1,g_2$, and set
\[
 E=\ker(g_1-1)\cap\ker(g_2-1),\qquad h_0=\dim_{\mathbf C}E.
\]
Every vector of $E$ is fixed by $w$, and (BI.34) shows that it
is fixed by all $s_j$.  Hence $h_0=\dim\tau^{H_0}$.
Moreover $E$ and its complex conjugate $\overline E$ are disjoint:
$A_1$ has eigenvalue $\lambda^{-1}$ on $E$ and eigenvalue
$\lambda$ on $\overline E$, and these are distinct because
$\lambda$ is nonreal.  Both subspaces lie in the real fixed space
of $w$, after complexification.  Thus
\[
                         2h_0\le4-\ell.                    \tag{BI.35}
\]

We now construct the degeneration and compute its graph module.
For every equal pair $(r_j,r_j)$ in (BI.33), bubble off a rational
base component containing those two branch points and one node.
Retain all $s_i,g_1,g_2$ on the central rational component.
Since $r_j^2=1$, the node inertia is trivial.  The local cover over
each leaf has monodromy $\langle r_j\rangle$ and is a disjoint
union of rational double covers.  The local cover over the central
component has monodromy $H_0$.  Gluing the unramified node fibres
produces the admissible $J$-cover determined by the tuple, and
smoothing the nodes recovers the given Hurwitz component.  All
base components are stable: each leaf has three special points,
and the central component has $\ell+2+q\ge3$.

Let $\Gamma$ be the dual graph of the source curve.  There is one
copy of the regular $J$-set of edges for each leaf, while its vertex
sets are $J/\langle r_j\rangle$ and $J/H_0$.  The graph is
connected because the original tuple, and hence
$H_0,r_1,\ldots,r_q$, generates $J$.  Its exact cellular sequence is
\[
 0\longrightarrow H_1(\Gamma,\mathbf C)
 \longrightarrow\mathbf C[J]^q
 \longrightarrow\mathbf C[J/H_0]\oplus
       \bigoplus_{j=1}^{q}\mathbf C[J/\langle r_j\rangle]
 \longrightarrow\mathbf C\longrightarrow0.
                                                               \tag{BI.36}
\]
Taking $\tau$-multiplicities is exact.  The regular representation
has multiplicity four, each reflection has a three-dimensional
fixed subspace, and $\tau$ has no invariant vector for $J$.
Therefore (BI.35)--(BI.36) give
\[
 \dim\operatorname{Hom}_J(\tau,H_1(\Gamma,\mathbf C))
       =4q-3q-h_0=q-h_0
       \ge\frac{10-\ell}{2}-\frac{4-\ell}{2}=3.
                                                               \tag{BI.37}
\]

For clarity, this graph dimension is the rank of one actual
monodromy logarithm.  Smooth all $q$ node parameters along a
one-parameter arc with positive integer orders.  The vanishing
cycles are disjoint.  The sum $N$ of their positively weighted
Picard--Lefschetz logarithms satisfies $N^2=0$; its rank on a
coefficient multiplicity equals the corresponding graph multiplicity.
One can see the equality directly from (BI.36): the independent
vanishing cycles form the edge-cycle space, and the monodromy
pairing on that space is the restriction of the positive diagonal
edge-length form.  It is nondegenerate, so no additional kernel
appears after taking the finite-group isotypic summand.  This is
the same semistable-curve graph calculation used in D-G34-03.
The smooth $\tau$-multiplicity has dimension six; since $N^2=0$,
its rank is at most three.  Equation (BI.37) forces
\[
                             \operatorname{rank}N=3.          \tag{BI.38}
\]

It remains to descend this conclusion to the original field $K$.
The intermediate Hurwitz component has exactly twelve branch points
and dimension nine.  The original period map factors through it
with rank nine, so the original family's classifying map to this
component is dominant.  Choose a subvariety of the original base
mapping dominantly and generically finitely onto the component;
one obtains it by choosing a closed point of the nonempty generic
fibre, whose residue field is a finite extension of the component's
function field, and taking its closure.  After shrinking a nonempty open and normalizing,
this gives a finite etale cover on which the original $K$-variation
is defined and whose chosen complex multiplicity is the intermediate
$\tau$-variation.

Choose the preceding one-parameter smoothing arc generically so
that its punctured part lies in this nonempty open.  After a finite
base change it lifts to that finite cover.  Its monodromy on the
original $K$-variation is therefore a positive power of the
intermediate monodromy.  The logarithm $N_K$ is $K$-linear,
has square zero, and its rank is three because the rank after the
chosen complex embedding is (BI.38).  This argument neither
identifies $L$ with $K$ nor asserts that an $L$-isotropic subspace
automatically descends.

The original polarization makes $N_K$ skew adjoint.  Hence
\[
 h(N_Kx,N_Ky)=-h(x,N_K^2y)=0.
\]
Its image is a three-dimensional totally isotropic $K$-subspace
of the six-dimensional original Hermitian space.  That Hermitian
form is split, contrary to the hypothesis.  This proves the
exclusion of all real reflection normal closures and leaves
exactly the nonreal rows (BI.31) from D-BI-11.
\end{proof}
